\chapter{Calcul algébrique et polynômes}
\section*{Exercice 2.1 -- Développer et contrôler}
Le développement donne $4x^2-12x+9-(x^2-1)=3x^2-12x+10$. En $0$, les deux formes donnent $10$ ; en $2$, elles donnent $-2$. Ces contrôles détectent certaines erreurs mais ne prouvent pas l'identité, déjà établie par développement.

\section*{Exercice 2.2 -- Domaine avant simplification}
Le domaine est $\R\setminus\{3\}$. Comme $x^2-9=(x-3)(x+3)$, l'expression vaut $x+3$ sur ce domaine. L'équation donnerait $x=3$, valeur interdite : elle n'a donc aucune solution. La simplification ne rétablit pas un point absent du domaine.

\section*{Exercice 2.3 -- Racine et facteur}
On regroupe $x^2(x-2)-(x-2)=(x-2)(x^2-1)=(x-2)(x-1)(x+1)$. Les racines sont donc $-1,1,2$, par la propriété du produit nul.

\section*{Exercice 2.4 -- Une identité générale}
En distribuant, la première somme est $\sum_{k=0}^{n-1}x^{n-k}a^k$ et la seconde $\sum_{k=0}^{n-1}x^{n-1-k}a^{k+1}$. En décalant l'indice de la seconde, les termes intermédiaires s'annulent. Seuls $x^n$ et $-a^n$ restent. La preuve vaut aussi pour $x=a$, sans division.

\section*{Exercice 2.5 -- Identifier des polynômes}
Leur différence $R$ a trois racines distinctes et un degré au plus $2$ si elle n'est pas nulle. Le théorème du nombre de racines d'un polynôme non nul rend cela impossible. Donc $R$ est le polynôme nul. Trois valeurs arbitraires ne suffiraient pas pour des fonctions générales.

\section*{Exercice 2.6 -- Lettres de rôles différents}
Développer le membre de gauche donne $x^2-(a+b)x+ab$. L'identité exige $a+b=5$ et $ab=6$. Les nombres $a,b$ sont donc les racines de $t^2-5t+6=(t-2)(t-3)$ : $(a,b)=(2,3)$ ou $(3,2)$. Ici $x$ parcourt tous les réels ; $a,b$ sont les paramètres cherchés.
